% Vérifié : renvois de feuillets contrôlés (tous exacts) ; «~veille~» de l'Ancre corrigé (onze jours) ; «~rappelé par un message~» ramené au texte ; «~81st~» et dysenterie précisés ; pince-nez (réparation + neuf, Imbert) ; ration d'avoine de 6 lb restreinte aux chevaux d'artillerie (OH 1917, p. 66) ; citation de Fernand-Laurent replacée dans son contexte (cantonnement) ; rumeur d'octobre sourcée (f. 42 r°) et abrégée (déjà dans l'intro générale) ; concerts payants attribués aux «~13th corps Archie Singers~» d'après le carnet ; Englebelmer vide : constat du carnet, non de la CWGC ; canonnade du 28 déc. : formulation atténuée. Style allégé (oct. 2026).
% Introduction du chapitre V (p. 105–127, ff. 52 v° – 63 v°).
Pendant ces huit semaines, Hippolyte Combes ne quitte guère Acheux-en-Amiénois, où il loge chez un tailleur, M.~Carluvant (lecture incertaine). Il n'est plus l'interprète d'une compagnie du train : il y accompagne en automobile un officier britannique chargé d'acheter de la paille et du foin pour l'armée. Ce service, commencé le 2~novembre 1916, onze jours avant la bataille de l'Ancre, devient une routine d'hiver. Brouillard et boue, puis grand gel à partir de la mi-janvier : les tournées se limitent à quelques villages entre l'Authie et l'Hallue. Seuls les heurts avec le jeune capitaine Williams, une lettre de Marguerite, les concerts du foyer britannique et quelques courses à Amiens ou à Doullens rompent la monotonie.

\subsection*{Le carnet}

Le 1\textsuperscript{er}~décembre, après un passage au bureau de la Mission française de la Fifth Army, à Toutencourt, puis sur un message trouvé à Acheux, Hippolyte Combes est «~présenté au bureau M.~F\textsuperscript{se} XIII\textsuperscript{e} corps et pris en charge~» ; il envisage alors une «~réintégration prochaine ds régiment d'origine~» (f.~52~v\textsuperscript{o}). Le lendemain arrive le capitaine Williams. Il remplace le lieutenant Curtis, qui rejoint la 152nd Company le 6 (ff.~52~v\textsuperscript{o}–53~v\textsuperscript{o}).

Les premières tournées vont loin : Franvillers et Buire-sur-Ancre le 1\textsuperscript{er} (f.~52~v\textsuperscript{o}), Marcelcave par Corbie le 3 (f.~53~r\textsuperscript{o}), Bayencourt, «~à 3 kilomètre de Gommecourt (Boches)~», le 5 (f.~53~v\textsuperscript{o}), Englebelmer, «~à peu près désert : plus d'habitants~», le 7 (f.~54~r\textsuperscript{o}), Saulty, sur la route d'Arras, le 9 (f.~54~v\textsuperscript{o}). Ce jour-là, il déjeune à Authie au mess de la «~no 3 Coy~» d'une division qu'il note «~81\textsuperscript{st}~» (sans doute la 31st, alors en ligne au XIII Corps). Les tournées se resserrent ensuite autour de quelques lieux : la ferme de Belle-Église, Raincheval, siège de l'«~Army P. Board~» (f.~56~r\textsuperscript{o}), Louvencourt, Beauquesne, Rubempré.

Le travail d'interprète consiste à chercher la paille «~de ferme en ferme~» (f.~62~r\textsuperscript{o}) et à traiter avec les cultivateurs. À Beauquesne, il négocie avec Mme~Capron la location d'une machine à battre ; le contrat est conclu le 5~janvier, jour où il déjeune avec les lieutenants Walsh et Morris (ff.~57~v\textsuperscript{o}, 60~r\textsuperscript{o}). Le 2~janvier, à la ferme Cuisinier de Louvencourt, il procède à l'«~expertise machine endommagée~» (f.~59~r\textsuperscript{o}). Le 13~décembre, il assiste à la mairie de Toutencourt à un «~meeting de refugiés d'Englebelmer~», en présence du major Levy (f.~55~r\textsuperscript{o}).

Les rapports avec Williams se dégradent vite : «~ce gosse de Williams~» (f.~54~v\textsuperscript{o}), «~1 cap. de 20 ans, présomptueux \& vain~» (f.~55~v\textsuperscript{o}). Le 24, à Louvencourt, chez Mme~Cuisinier, le capitaine lâche un «~Bloody France etc.~!~», et Hippolyte Combes hésite «~à le reporter à French Mission~» (f.~56~v\textsuperscript{o}). Le 30, à Senlis, après un nouveau «~mot très regrettable~», une «~explication très calme, mais très ferme~» lui vaut des excuses (f.~58~v\textsuperscript{o}).

La guerre reste proche. Le 6~janvier, les routes sont encombrées de «~deplacement de troupe, convois d'artillerie, de A.S.C.~» ; le 9, près d'Hédauville, sous une «~très violente canonnade~», il retrouve la 150th Company du train de la 18th Division (f.~60~v\textsuperscript{o}) ; le 10, le canon tonne «~s/ ligne de Miraumt~», et il retrouve «~enfin mon vieux Morris~» (f.~61~r\textsuperscript{o}). Le 13, au bureau de la Mission du «~V\textsuperscript{e}~corps~», il voit le capitaine Huguet : «~mon départ fixé definit au 26~ct~» (f.~61~v\textsuperscript{o}).

L'humeur est d'abord à l'ennui : «~Heures mornes en auto : pu ruminer \& observer à loisir~» (f.~53~r\textsuperscript{o} ; cf. \textit{Feuillets de Route}, chap.~\textsc{x} et~\textsc{xii}, annexe, p.~\pageref{annexe:doc6}) ; «~Jours se suivent, monotones \& ternes~» (f.~55~v\textsuperscript{o}). Le 22~décembre : «~Lettre de M\textsuperscript{te} m'a bp deplu~» (f.~56~v\textsuperscript{o}). Quelques courses plus lointaines le distraient : Amiens le 3~janvier, avec une longue visite de la cathédrale et, en route, une conversation avec le maire de Montigny, M.~Paul Dubois (f.~59~v\textsuperscript{o}) ; Doullens le 15 (f.~61~v\textsuperscript{o}) ; Auxi-le-Château le 17, où M.~Renard lui montre un «~coffret en paille acheté en 1858 au bagne de Toulon~» (f.~62~r\textsuperscript{o} ; cf. \textit{Feuillets de Route}, chap.~\textsc{ii} : «~Six Semaines (de bagne) à Toulon~», annexe, p.~\pageref{annexe:doc6}).

La solde, de 78~F~75 ou de 84~F par quinzaine, est versée par le bureau de corps de la Mission (ff.~53~r\textsuperscript{o}, 56~r\textsuperscript{o}, 60~r\textsuperscript{o}, 62~v\textsuperscript{o}). Le mess coûte 88~F en décembre (total inscrit par Hippolyte Combes ; les termes notés font 78~F) et 79~F~10 en janvier (ff.~58~v\textsuperscript{o} et 63~v\textsuperscript{o}). Le 24~janvier, les vitres de sa chambre sont couvertes de glace ; le 25, «~Ts préparatifs terminés~», il est prêt à partir (f.~63~v\textsuperscript{o}).

\subsection*{Le contexte}

Après la bataille de l'Ancre (13--18~novembre 1916), le front de la Fifth Army est immobile. En décembre, écrit Haig, l'effort porte sur «~l'amélioration de nos tranchées et des routes et communications~»\footnote{J.~H.~Boraston (éd.), \textit{Sir Douglas Haig's Despatches (December 1915 – April 1919)}, Londres et Toronto, J.~M.~Dent \& Sons, 1919, dépêche du 31~mai 1917, p.~65--66. Citations anglaises traduites par l'éditeur, sauf quand elles sont données en anglais.} ; l'état du champ de bataille, surtout dans la vallée de l'Ancre, ne fut «~probablement jamais dépassé~» sur le front occidental\footnote{W.~Miles, \textit{Military Operations, France and Belgium, 1916}, vol.~II, Londres, Macmillan, 1938 (histoire officielle britannique), p.~536--537.}. Au début de janvier, le secteur que longent ses tournées, au nord de l'Ancre, est tenu par le XIII Corps du lieutenant-général Congreve (7th, 3rd et 31st Divisions) ; en face, Gommecourt reste un saillant allemand jusqu'au 27~février 1917\footnote{C.~Falls, \textit{Military Operations, France and Belgium, 1917}, vol.~I, Londres, Macmillan, 1940 (ci-après \textit{OH 1917}, I), p.~66, 95--101 et 115. Bayencourt est en réalité à un peu plus de 5~km de Gommecourt (calcul d'après la position des mairies).}. La canonnade du 28~décembre vers Hébuterne ne correspond à aucune opération signalée par l'histoire officielle. Celles des 9 et 10~janvier préparent la prise du Triangle, à l'est de Beaumont-Hamel, et de Munich Trench (10--11~janvier) : le tir lourd fut «~maintenu pendant les deux jours précédents, avec de rares interruptions, sur tout le front de la Fifth Army~»\footnote{\textit{OH 1917}, I, p.~67--69.}. Suivent les relèves : le V Corps du lieutenant-général Fanshawe reprend le secteur du XIII, la 32nd Division relevant la 3rd le 9 et la 19th la 31st le 11\footnote{\textit{Ibid.}, p.~66 et 69.} ; les convois du 6 en sont vraisemblablement le prélude. Pour la fin de janvier, l'histoire officielle ne signale «~aucun événement de quelque importance~» ; le «~violent bombardemt~» du 25 reste inexpliqué\footnote{\textit{Ibid.}, p.~69.}.

Le temps noté par Hippolyte Combes concorde avec ces sources : douceur du début de janvier, gel à partir du 14, «~gelé extr.~+ fort~» le 25, jour d'un minimum de 15~°F (environ $-$9{,}4~°C). Ce gel dura un mois et ruina les routes peu empierrées de l'arrière\footnote{\textit{Ibid.}, p.~66 et 69.}. Commencée «~brusquement~» vers le 16~janvier, la vague de froid toucha toute la France jusque vers le 16~février\footnote{G.~Barbé, «~L'hiver 1916-1917~», \textit{Journal d'agriculture pratique}, nouv. série, t.~30, 1917, p.~181--182.}.

Le passage du bureau «~XIII\textsuperscript{e}~corps~» au bureau «~V\textsuperscript{e}~corps~» de la Mission suit les états-majors : le XIII Corps avait remplacé le V Corps à Acheux le 29~novembre 1916 (voir le chapitre~\textsc{iv}), et les deux quartiers généraux échangent de nouveau leurs places le 12~janvier 1917\footnote{C.~Baker, \textit{The Long, Long Trail}, notices «~The British V Corps in France and Flanders~» et «~The British XIII Corps in France and Flanders~», \url{https://www.longlongtrail.co.uk} (consultées le 18~octobre 2026).}. Le retour au régiment envisagé le 1\textsuperscript{er}~décembre fait écho à la rumeur du 22~octobre sur le renvoi des interprètes R.A.T. (f.~42~r\textsuperscript{o} ; voir le chapitre~\textsc{iv}). Le carnet ne dit pas si Hippolyte Combes signala l'incident du 24~décembre. Un autre interprète reconnaissait que des «~heurts désagréables~», «~dont les interprètes sont toujours les victimes~», se produisaient trop souvent entre Britanniques et Français, à propos du cantonnement il est vrai\footnote{Fernand-Laurent, \textit{Chez nos alliés britanniques. Notes et souvenirs d'un interprète}, préface de F.~de Tessan, Paris, Boivin, 1917, p.~231--232.}.

Le \textit{purchase board} de chaque armée achetait la paille, complément de la ration de fourrage et couchage des troupes, et il «~administers a number of baling and threshing machines~»\footnote{F.~K.~Puckle, \textit{The Army Service Corps of the British Army and the Organization of the Transport and Transportation at the Front in France}, Washington, Government Printing Office, 1917, p.~28.} : d'où le battage d'Arquèves et la batteuse de Beauquesne. Pendant la crise des transports de l'hiver, la ration d'avoine des chevaux de l'artillerie était tombée à six livres\footnote{\textit{OH 1917}, I, p.~65--66.} ; pour préparer Arras, on recommandera encore de «~vider rapidement les granges en aidant les fermiers à battre leur grain et en achetant la paille pour l'armée~»\footnote{\textit{Ibid.}, p.~190.}. Englebelmer, poste d'ambulances exposé à des bombardements jusqu'à l'automne de 1916\footnote{Commonwealth War Graves Commission, notice «~Englebelmer Communal Cemetery Extension~», \url{https://www.cwgc.org/visit-us/find-cemeteries-memorials/cemetery-details/4003798/englebelmer-communal-cemetery-extension/} (consultée le 17~octobre 2026).}, n'avait plus d'habitants, comme le constate Hippolyte Combes. L'objet de la réunion de Toutencourt, non précisé, touchait probablement les récoltes ou les fourrages laissés par les réfugiés. Pour toute réquisition ou réclamation, l'officier britannique devait d'abord trouver «~the Maire or the Garde-Champêtre~» : d'où les visites aux maires\footnote{«~Claims~», par «~A Divisional Claims Officer~», \textit{Blackwood's Magazine}, vol.~201, n\textsuperscript{o}~1220, juin 1917, p.~935--939.}.

La division d'Hippolyte Combes est loin. Relevée par la 61st Division, la 18th est partie «~for Christmas rest in the Abbeville area~» ; le 7th Bedfordshire cantonne à Domvast du 14~décembre au 10~janvier et y exécute, le 27~décembre, une attaque de démonstration devant «~the staff of the 5th Army \& Corps~»\footnote{G.~H.~F.~Nichols, \textit{The 18th Division in the Great War}, Édimbourg et Londres, Blackwood, 1922, p.~136 et 139 ; journal de marche du 7th Bedfordshire, décembre 1916, Bedfordshire Archives, X550/8/1/1612, \url{https://bedsarchivescat.bedford.gov.uk/Details/archive/110153355}.}. Le 15~janvier, le major-général Maxse cède le commandement au major-général R.~P.~Lee ; le 16, la division revient sur l'Ancre, au sud de la rivière, dans le II Corps\footnote{A.~F.~Becke, \textit{Order of Battle of Divisions}, part~3A, Londres, HMSO, 1938, p.~79 ; Nichols, \textit{op.~cit.}, p.~136--139 ; \textit{OH 1917}, I, p.~69--70.}. La 150th Company retrouvée le 9 près d'Hédauville et le «~vieux Morris~», vraisemblablement l'un de ses officiers, annoncent ce retour. Les capitaines Williams et Huguet ne sont pas identifiés.

Le foyer de l'Y.M.C.A. accueillait le soir des concerts d'amateurs\footnote{Fernand-Laurent, \textit{op.~cit.}, p.~90--95.}. Hippolyte Combes paie pourtant 30~centimes pour ceux des 28~décembre, 4 et 6~janvier, l'entrée étant en principe gratuite dans les huttes ordinaires\footnote{A.~K.~Yapp, \textit{The Romance of the Red Triangle}, Londres, Hodder and Stoughton, 1918, chap.~\textsc{iv}.} ; ces concerts étaient donnés par une troupe de corps d'armée, les «~13th corps Archie Singers~», le 6 reprenant le «~Programme des 2 fois precédentes~» (ff.~58~r\textsuperscript{o}, 60~r\textsuperscript{o} et 60~v\textsuperscript{o}).

Le départ du 26~janvier relève du régime des permissions en vigueur depuis le 1\textsuperscript{er}~octobre 1916 (voir le chapitre~\textsc{vi}). Le chapitre suivant le conduit, par la gare de Varennes et une longue attente à Gézaincourt, à Paris, à Toulouse et à Millau, auprès de Marguerite.
